%-------------------------------------------------------------------------------
%                            BAB V
%               		KESIMPULAN DAN SARAN
%-------------------------------------------------------------------------------
\chapter{KESIMPULAN DAN SARAN}

\section{Kesimpulan}Penelitian ini berhasil mengembangkan dan mengevaluasi metode analisis LIBS berbasis model DL yang dilatih sepenuhnya pada spektrum emisi sintetis. Pendekatan ini bertujuan untuk mengatasi ketergantungan pada kalibrasi eksperimental yang ekstensif. Validasi dataset sintetis menunjukkan kemampuannya mereplikasi kompleksitas spektral, termasuk fenomena interferensi, serta akurasi fisik yang tinggi dibandingkan dengan data referensi NIST, menegaskan fondasi yang kokoh untuk pelatihan model.Evaluasi komprehensif terhadap arsitektur model \textit{Informer} mengungkapkan temuan penting mengenai dampak kompleksitas model terhadap kinerja. Secara spesifik, model \textit{Informer} dengan \textbf{2-layer} menunjukkan kinerja yang secara signifikan lebih unggul dibandingkan dengan model \textbf{3-layer} pada metrik-metrik kunci seperti F1-Score (Macro), Precision (Macro), Recall (Macro), dan Subset Accuracy. Model 2-layer mencapai F1-Score (Macro) sebesar 0.5785, jauh melampaui 0.3274 yang dicapai oleh model 3-layer. Peningkatan ini mengindikasikan bahwa, dalam konteks dataset dan tugas identifikasi multi-elemen ini, arsitektur yang lebih sederhana mampu mempelajari representasi yang lebih efektif dan menggeneralisasi lebih baik. Hal ini menyoroti pentingnya optimisasi arsitektur model dan menunjukkan bahwa penambahan kompleksitas tidak selalu berkorelasi positif dengan peningkatan performa, terutama jika tidak diimbangi dengan data pelatihan yang memadai atau strategi regularisasi yang tepat.Keberhasilan model 2-layer dalam mencapai akurasi prediksi multi-elemen yang tinggi tanpa kalibrasi eksperimental, didukung oleh pelatihan eksklusif pada spektrum sintetis, secara langsung menjawab rumusan masalah pertama penelitian. Meskipun model menunjukkan kemampuan \textit{recall} yang tinggi, tantangan pada \textit{precision} untuk beberapa elemen minor tetap ada, mengindikasikan area untuk perbaikan lebih lanjut. Efisiensi komputasi model \textit{Informer}, yang merupakan salah satu tujuan penelitian, juga terbukti melalui kemampuannya memproses spektrum resolusi tinggi.Secara keseluruhan, penelitian ini menegaskan potensi besar pendekatan berbasis DL dengan data sintetis untuk analisis spektral LIBS. Ini membuka jalan bagi pengembangan metode analisis unsur yang lebih efisien, hemat biaya, dan berpotensi sepenuhnya bebas kalibrasi, merevolusi aplikasi LIBS di berbagai bidang.

\section{Saran}
Berdasarkan temuan dan keterbatasan yang telah diidentifikasi dalam penelitian ini, beberapa arah untuk penelitian selanjutnya dapat direkomendasikan untuk membangun di atas fondasi yang telah diletakkan:
\begin{enumerate}
    \item \textbf{Validasi dengan Sampel Dunia Nyata:} Mengingat keterbatasan model yang saat ini hanya diuji pada data sintetis, langkah paling krusial berikutnya adalah melakukan validasi ekstensif menggunakan data spektral eksperimental. Pengujian terhadap sampel nyata dengan beragam komposisi matriks akan secara langsung menguji ketahanan dan kemampuan generalisasi model dalam menghadapi efek matriks yang tidak terduga, yang merupakan salah satu keterbatasan utama dari pendekatan berbasis simulasi.
    \item \textbf{Perluasan Database dan Kondisi Simulasi:} Untuk mengatasi keterbatasan cakupan elemen dan asumsi LTE, penelitian selanjutnya disarankan untuk:
    \begin{itemize}
        \item Memperluas dataset sintetis dengan menyertakan lebih banyak elemen, terutama elemen-elemen penting secara industri atau lingkungan yang belum tercakup.
        \item Mengembangkan model simulasi yang lebih canggih dengan mempertimbangkan kondisi non-LTE, yang dapat meningkatkan realisme spektrum yang dihasilkan dan akurasi model pada fase plasma yang lebih dinamis.
    \end{itemize}
    \item \textbf{Pengembangan Analisis Kuantitatif:} Saat ini, model berfokus pada identifikasi kualitatif. Sebuah langkah maju yang signifikan adalah mengembangkan arsitektur yang mampu melakukan analisis kuantitatif untuk memprediksi konsentrasi absolut atau relatif dari setiap elemen. Ini akan secara drastis meningkatkan nilai praktis dari metode yang diusulkan.
    \item \textbf{Eksplorasi Arsitektur dan Interpretabilitas:} Meskipun \textit{Informer} terbukti efisien, penelitian di masa depan dapat mengeksplorasi arsitektur alternatif seperti model hibrida (CNN-Transformer) yang mungkin lebih baik dalam menangkap fitur spektral lokal dan dependensi global. Selain itu, penerapan teknik interpretabilitas (misalnya, SHAP atau LIME) sangat direkomendasikan untuk membuka "kotak hitam" model, memberikan pemahaman fisis tentang fitur spektral apa yang menjadi dasar prediksi, dan pada akhirnya, meningkatkan kepercayaan terhadap metode ini.
\end{enumerate}
